\chapter{Conclusions}
\label{ch:con}
%
Four PDE problems were solved with four different deep learning methods,
all implemented from scratch in PyTorch: an inverse PINN with a two-stage
curriculum for coefficient identification from noisy data (Problem A,
relative $L^2$ error $1.18\times10^{-1}$ for $k$ and $5.1\times10^{-3}$ for
$u$), the Deep Ritz method for a forward problem with a discontinuous
coefficient (Problem B, $3.2\times10^{-2}$), a Fourier Neural Operator with
a hand-written truncated DFT for supervised operator learning (Problem C,
$7.2\times10^{-3}$), and a physics-informed DeepONet for semi-supervised
operator learning (Problem D, $4.1\times10^{-2}$).

Beyond the individual results, three observations recur across the
problems. First, hard constraints beat penalties whenever they are
available: building boundary conditions, positivity, or known structure
directly into the ansatz removed loss weights and improved stability in
every problem where we used it. Second, the choice of loss formulation
matters more than network size for difficult coefficients: for the
discontinuous permeability of Problem B, moving from a strong-form residual
to an energy formulation was the difference between failure and a $3\%$
solution, while wider networks changed little. Third, curricula help
non-convex training: both in Problem A (fit the state first, then identify
the coefficient) and in Problem D (fit the data first, then add the
residual), staged training escaped plateaus that joint training from a
random initialization did not; in Problem D the physics phase reduced the
error of the supervised plateau by a factor of $2.5$ by exploiting the
$1800$ unlabeled initial conditions.

The main limitations, and natural directions for future work, follow from
the same experiments. The accuracy of the recovered coefficient in Problem
A is bounded by the noise level at the point where the flux degenerates;
ensembling or a Bayesian treatment of the state fit could quantify that
uncertainty instead of hiding it. The smooth tanh ansatz in Problem B
cannot represent gradient kinks at material interfaces; enriching the
ansatz with interface-aware features would address the dominant error
bands. Finally, the FNO in Problem C shows a mild train--test gap that more
training data or stronger regularization would likely close.
